%==================================================
%----- 2.7 Mínima reação — conexão dinâmica
%==================================================
\section{Mínima reação}
Em ausência de torque externo sobre a base, o momento angular total é constante.
Escrevendo o momento angular da base como:
\begin{equation}
\vec h_b
=
\mathbf I_b(\vec q)\,\vec\omega_b
+\mathbf H_b(\vec q)\,\dot{\vec q},
\label{eq:hb_def}
\end{equation}
onde $\mathbf I_b(\vec q)$ é a inércia efetiva da base e
$\mathbf H_b(\vec q)$ é a matriz de acoplamento base--manipulador,
a condição de mínima reação busca manter
$\dot{\vec h}_b=\vec 0$
durante o movimento do manipulador.
Ao derivar \eqref{eq:hb_def}, obtém-se:
\begin{equation}
\mathbf H_b(\vec q)\,\ddot{\vec q}
+\dot{\mathbf H}_b(\vec q,\dot{\vec q})\,\dot{\vec q}
+\mathbf I_b(\vec q)\,\dot{\vec\omega}_b
=
\vec 0.
\label{eq:hb_dot_zero}
\end{equation}

Se a atitude da base é regulada para $\vec\omega_b\!\approx\!\vec 0$,
adota-se a seguinte condição:
\begin{equation}
\mathbf H_b(\vec q)\,\ddot{\vec q}
+\dot{\mathbf H}_b(\vec q,\dot{\vec q})\,\dot{\vec q}
=
\vec 0.
\label{eq:min_react_cond}
\end{equation}

\subsection{Projeção no espaço nulo}
No nível de velocidade, pode-se expressar como:
\begin{equation}
\mathbf H_b(\vec q)\,\dot{\vec q}
=
\vec 0,
\qquad
\dot{\vec q}
=
\mathbf N_b(\vec q)\,\vec v,
\label{eq:vel_null}
\end{equation}
com o projetor ortogonal
$\mathbf N_b(\vec q)=\mathbf I-\mathbf H_b^{+}(\vec q)\,\mathbf H_b(\vec q)$,
sendo $\mathbf H_b^{+}$ a pseudoinversa de Moore--Penrose.
Assim, a referência de velocidade $\dot{\vec q}$ não possar injetar momento na base.

No nível de aceleração, para acompanhar uma tarefa principal com aceleração desejada
$\ddot{\vec q}_{\text{t}}$,
usa-se:
\begin{equation}
\ddot{\vec q}
=
\ddot{\vec q}_{\text{t}}
-\mathbf H_b^{+}(\vec q)\,\dot{\mathbf H}_b(\vec q,\dot{\vec q})\,\dot{\vec q}
+\mathbf N_b(\vec q)\,\vec w,
\label{eq:acc_null}
\end{equation}
onde o segundo termo compensa a parcela convectiva de
\eqref{eq:min_react_cond}
e $\mathbf N_b\,\vec w$
utiliza os graus de liberdade residuais para objetivos secundários, como evitar juntas-limite, suavizar torque, etc.

\subsection{Compatibilização com vínculos holonômicos}
Se houver restrições $\Phi(\vec q,t)=\vec 0$,
tratadas como em \eqref{eq:vellevel},
e \eqref{eq:acclevel},
ou pela estabilização de Baumgarte
\eqref{eq:acclevel_baumgarte},
combina-se \eqref{eq:min_react_cond}
com o vínculo de aceleração resultando no sistema linear:
\begin{equation}
\begin{bmatrix}
\mathbf H_b(\vec q)\\[2pt]
\mathbf J(\vec q,t)
\end{bmatrix}
\ddot{\vec q}
\;=\;
\begin{bmatrix}
-\dot{\mathbf H}_b(\vec q,\dot{\vec q})\,\dot{\vec q}\\[2pt]
-\dot{\mathbf J}(\vec q,\dot{\vec q},t)\,\dot{\vec q}
-\dfrac{\partial^{2}\Phi}{\partial t^{2}}
-2\alpha\!\left(\mathbf J(\vec q,t)\,\dot{\vec q}+\dfrac{\partial \Phi}{\partial t}\right)
-\beta^{2}\Phi(\vec q,t)
\end{bmatrix},
\label{eq:stack_min_react_baumgarte}
\end{equation}
que pode ser resolvido por mínimos quadrados ponderados. De forma alternativa, pode-se aplicar apenas a projeção
\eqref{eq:vel_null}
ao gerador de referências e manter o fechamento de malha nas juntas.

\subsection{Observações}
(i) $\dot{\mathbf H}_b$ é calculado de modo consistente com a cinemática adotada;  
(ii) $\mathbf H_b^{+}$ foi regularizado em perda de posto (pseudoinversa amortecida);  
(iii) Como a base também possui controle, foi considerado o termo
$\mathbf I_b(\vec q)\,\dot{\vec\omega}_b$
de \eqref{eq:hb_dot_zero}
para não conflitar com o regulador da base.